\chapter{Introduction}
\label{ch:introduction}

\section{Background}

Chiang Mai, located in the mountainous basin of northern Thailand,
experiences one of the most acute seasonal air-pollution episodes in
Southeast Asia. Between February and April each year, the combination of
widespread agricultural burning, forest fires, and stagnant boundary-layer
meteorology drives PM2.5 concentrations well above both the World Health
Organization's 24-hour guideline of 15~\textmu g/m\textsuperscript{3}
(\cite{who2021}) and Thailand's national standard of
50~\textmu g/m\textsuperscript{3}. Concentrations above
150~\textmu g/m\textsuperscript{3} are routine during the peak burning weeks,
with daily peaks occasionally exceeding 300~\textmu g/m\textsuperscript{3}.

The public-health impact is substantial. Elevated PM2.5 is causally linked
to respiratory disease, cardiovascular events, and premature mortality, with
particularly acute risk for children, the elderly, pregnant women, and
outdoor workers. Against this backdrop, Thailand's Pollution Control
Department (PCD, \cite{pcd}) publishes measured PM2.5 values through the
air4thai portal, but the portal reports observed concentrations rather than
forecasts. Residents are therefore informed that conditions are already
hazardous; they cannot plan several days ahead around predicted hazardous
episodes, nor is there a systematic, publicly accessible multi-horizon
forecast service for Chiang Mai.

This gap --- between rich, available measurement data and the absence of a
production-grade forecasting service --- defines the opportunity that the
present project addresses.

\section{Problem Statement}

Despite the severity of PM2.5 pollution in Chiang Mai, no publicly
accessible automated forecasting service provides residents with reliable
multi-day PM2.5 predictions or advance hazard alerts. Concretely, the
problem addressed in this project can be stated as follows:

\begin{quote}
Given fifteen years of historical PM2.5 measurements at a Chiang Mai
monitoring station, co-located meteorological data, and regional satellite
fire-hotspot observations, build a production-grade machine-learning system
that (a) forecasts PM2.5 one, three, and seven days ahead, (b) issues binary
hazard alerts when forecast PM2.5 exceeds 50~\textmu g/m\textsuperscript{3},
and (c) operates under a continuously monitored, self-retraining MLOps
pipeline deployable within a student-scale cloud budget.
\end{quote}

The problem is not merely a modelling exercise: reaching a grade of
\emph{production-grade} requires addressing data engineering
(heterogeneous-source ingestion, deduplication, schema stability),
machine-learning discipline (chronological splits, champion/secondary
governance, interpretability), MLOps rigour (drift monitoring with
seasonal awareness, automated retraining, infrastructure as code), and
operational economics (sub-dollar-per-day cost to justify realistic
deployment).

\section{Research Questions}

This project is organised around four concrete research questions that
motivate its design and evaluation:

\begin{enumerate}
    \item \textbf{RQ1.} Can a machine-learning model trained on daily
          PM2.5, weather, and fire-hotspot features forecast PM2.5 for
          Chiang Mai at t+1, t+3, and t+7 horizons with meaningfully lower
          error than a persistence baseline?

    \item \textbf{RQ2.} Given that both gradient-boosted tree ensembles
          and recurrent deep sequence models are widely used for tabular
          time-series forecasting, which family performs better on this
          specific problem, and is there measurable benefit from retaining
          more than one model family in production?

    \item \textbf{RQ3.} What is the measurable contribution of NASA FIRMS
          satellite fire-hotspot features to forecast accuracy, and does
          that contribution vary across forecast horizons?

    \item \textbf{RQ4.} Can a complete, self-retraining MLOps system ---
          including automated ingestion, drift monitoring, champion/challenger
          validation, and cloud deployment --- be operated within a
          student-scale budget while remaining rubric-compliant and
          reproducible end-to-end?
\end{enumerate}

\section{Objectives}

The above research questions translate into the following project
objectives. This project --- the
\textbf{ChiangMai Air Quality Alert System (CAAS)} --- aims to:

\begin{enumerate}
    \item \textbf{Build an end-to-end data-engineering pipeline} that
          ingests, consolidates, and feature-engineers PM2.5, weather, and
          fire data from three independent public sources, producing a
          clean 45-feature daily matrix spanning 2011--2025.

    \item \textbf{Develop, tune, and compare multiple forecast model
          families} --- LightGBM, XGBoost, and LSTM --- under a shared
          chronological evaluation protocol at t+1, t+3, and t+7 horizons.

    \item \textbf{Issue binary hazard alerts} at a physically motivated
          threshold (PM2.5 $>$ 50~\textmu g/m\textsuperscript{3}) and
          empirically evaluate whether this default threshold is optimal
          through precision--recall analysis.

    \item \textbf{Operationalise the winning model(s)} behind a FastAPI
          inference service with MLflow experiment tracking, a Streamlit
          monitoring dashboard, Evidently drift detection, and GitHub Actions
          CI/CD orchestration.

    \item \textbf{Provision all cloud infrastructure through Terraform}
          such that the full stack can be deployed to a fresh AWS account
          with a single \texttt{terraform apply}.

    \item \textbf{Interpret the resulting model} with SHAP analysis to
          demonstrate that the dominant predictive signal shifts in a
          physically interpretable way with forecast horizon.
\end{enumerate}

\section{Scope and Contributions}

The scope of this project is deliberately bounded to keep it deliverable
within a single capstone semester. CAAS targets a single PCD monitoring
station in Chiang Mai (station codes 35T/36T), operates at daily temporal
resolution, and uses historical weather reanalysis rather than live weather
forecast inputs. Multi-station extension, hourly resolution, and forecast
weather are discussed as future work in Chapter~\ref{ch:conclusion}.

Within this scope, the key contributions of this project are:

\begin{itemize}
    \item A 45-feature dataset spanning 5,257 days (2011--2025), merging
          three heterogeneous sources: PCD PM2.5 measurements, Open-Meteo
          weather reanalysis, and NASA FIRMS fire hotspot counts.

    \item A LightGBM champion (Optuna-tuned) and XGBoost strong-secondary,
          both operationalised behind the same FastAPI service. The champion
          achieves MAE of 5.12, 6.77, and 8.25~\textmu g/m\textsuperscript{3}
          at t+1, t+3, and t+7 respectively, with
          AUROC~$\geq$~0.95 across all horizons.

    \item A SHAP-based interpretability analysis demonstrating that the
          dominant predictive signal shifts from recent PM2.5 lags at
          short horizons to fire accumulation and seasonal calendar
          position at longer horizons.

    \item An ablation study quantifying the contribution of FIRMS fire
          data: removing the six fire-derived features degrades MAE by
          4.3--6.8\% across horizons.

    \item A complete MLOps stack --- FastAPI inference service, Streamlit
          dashboard, Evidently drift monitoring with seasonal-aware
          retrain policy, Terraform-provisioned AWS infrastructure,
          and GitHub Actions CI/CD --- operating at approximately
          \$0.55/day.
\end{itemize}

\section{Organisation of the Report}

The remainder of this report is organised as follows.
Chapter~\ref{ch:literature-review} reviews prior work on PM2.5
forecasting, gradient-boosted versus deep sequence models for tabular
time-series, satellite fire observations, and MLOps patterns for
environmental ML systems.
Chapter~\ref{ch:data-engineering} describes the three data sources, the
ingestion pipeline, the 45-feature engineering, and the S3-based storage
and versioning layout.
Chapter~\ref{ch:ml-development} details model selection, hyperparameter
tuning, evaluation results, ablation studies, alert-threshold analysis,
and SHAP interpretability.
Chapter~\ref{ch:mlops} describes the MLOps architecture ---
MLflow tracking, FastAPI serving, CI/CD pipelines, Evidently drift
monitoring, and Terraform infrastructure as code.
Chapter~\ref{ch:performance} presents performance and scalability
analysis, inference latency, and a cost breakdown at AWS on-demand and
AWS Learner Lab pricing.
Chapter~\ref{ch:conclusion} concludes the report, summarises the key
findings, discusses limitations, and outlines recommendations and
future research directions. The References section follows, and the
Appendix provides a complete step-by-step reproduction guide and full
metric tables.
